\documentclass[10pt,a4paper]{amsart}
\usepackage{amsmath,amssymb,amsthm,graphicx}
% Retained minimal preamble for the Stacks Project chapter varieties.tex.
% Verbatim excerpt of the chapter preamble: Xy-pic setup (native preamble lines 14-18), the native theorem
% declarations with the mathematical macros (native preamble lines 154-217), and the project own no-op
% definitions of the reference/slogan/history display environments (native wrapper.tex lines 1-3) that the
% counted statement uses for its slogan. Copied mechanically; nothing added, removed or redefined.
% For commutative diagrams we use Xy-pic
\usepackage[all]{xy}

% We use 2cell for 2-commutative diagrams.
\xyoption{2cell}
% Theorem environments.
%
\theoremstyle{plain}
\newtheorem{theorem}[subsection]{Theorem}
\newtheorem{proposition}[subsection]{Proposition}
\newtheorem{lemma}[subsection]{Lemma}

\theoremstyle{definition}
\newtheorem{definition}[subsection]{Definition}
\newtheorem{example}[subsection]{Example}
\newtheorem{exercise}[subsection]{Exercise}
\newtheorem{situation}[subsection]{Situation}

\theoremstyle{remark}
\newtheorem{remark}[subsection]{Remark}
\newtheorem{remarks}[subsection]{Remarks}

\numberwithin{equation}{subsection}

% Macros
%
\def\lim{\mathop{\mathrm{lim}}\nolimits}
\def\colim{\mathop{\mathrm{colim}}\nolimits}
\def\Spec{\mathop{\mathrm{Spec}}}
\def\Hom{\mathop{\mathrm{Hom}}\nolimits}
\def\Ext{\mathop{\mathrm{Ext}}\nolimits}
\def\SheafHom{\mathop{\mathcal{H}\!\mathit{om}}\nolimits}
\def\SheafExt{\mathop{\mathcal{E}\!\mathit{xt}}\nolimits}
\def\Sch{\mathit{Sch}}
\def\Mor{\mathop{\mathrm{Mor}}\nolimits}
\def\Ob{\mathop{\mathrm{Ob}}\nolimits}
\def\Sh{\mathop{\mathit{Sh}}\nolimits}
\def\NL{\mathop{N\!L}\nolimits}
\def\CH{\mathop{\mathrm{CH}}\nolimits}
\def\proetale{{pro\text{-}\acute{e}tale}}
\def\etale{{\acute{e}tale}}
\def\QCoh{\mathit{QCoh}}
\def\Ker{\mathop{\mathrm{Ker}}}
\def\Im{\mathop{\mathrm{Im}}}
\def\Coker{\mathop{\mathrm{Coker}}}
\def\Coim{\mathop{\mathrm{Coim}}}

% Boxtimes
%
\DeclareMathSymbol{\boxtimes}{\mathbin}{AMSa}{"02}

%
% Macros for moduli stacks/spaces
%
\def\QCohstack{\mathcal{QC}\!\mathit{oh}}
\def\Cohstack{\mathcal{C}\!\mathit{oh}}
\def\Spacesstack{\mathcal{S}\!\mathit{paces}}
\def\Quotfunctor{\mathrm{Quot}}
\def\Hilbfunctor{\mathrm{Hilb}}
\def\Curvesstack{\mathcal{C}\!\mathit{urves}}
\def\Polarizedstack{\mathcal{P}\!\mathit{olarized}}
\def\Complexesstack{\mathcal{C}\!\mathit{omplexes}}
% \Pic is the operator that assigns to X its picard group, usage \Pic(X)
% \Picardstack_{X/B} denotes the Picard stack of X over B
% \Picardfunctor_{X/B} denotes the Picard functor of X over B
\def\Pic{\mathop{\mathrm{Pic}}\nolimits}
\def\Picardstack{\mathcal{P}\!\mathit{ic}}
\def\Picardfunctor{\mathrm{Pic}}
\def\Deformationcategory{\mathcal{D}\!\mathit{ef}}
\newenvironment{reference}{}{}
\newenvironment{slogan}{}{}
\newenvironment{history}{}{}
\usepackage{cleveref}
\crefname{originaltheorem}{original theorem}{original theorems}
\Crefname{originaltheorem}{Original theorem}{Original theorems}
\makeatletter
\newlabel{algebra-theorem-nullstellensatz}{{00FV}{source}{}{sourcepointer-1}{}}
\newlabel{algebra-theorem-nullstellensatz@cref}{{[originaltheorem][0][]00FV}{[1][0][]source}{}{}{}}
\makeatother
\crefname{originallemma}{original lemma}{original lemmas}
\Crefname{originallemma}{Original lemma}{Original lemmas}
\makeatletter
\newlabel{morphisms-lemma-base-change-finite-type}{{01T4}{source}{}{sourcepointer-2}{}}
\newlabel{morphisms-lemma-base-change-finite-type@cref}{{[originallemma][0][]01T4}{[1][0][]source}{}{}{}}
\makeatother
\crefname{originallemma}{original lemma}{original lemmas}
\Crefname{originallemma}{Original lemma}{Original lemmas}
\makeatletter
\newlabel{morphisms-lemma-composition-finite-type}{{01T3}{source}{}{sourcepointer-3}{}}
\newlabel{morphisms-lemma-composition-finite-type@cref}{{[originallemma][0][]01T3}{[1][0][]source}{}{}{}}
\makeatother
\crefname{originallemma}{original lemma}{original lemmas}
\Crefname{originallemma}{Original lemma}{Original lemmas}
\makeatletter
\newlabel{properties-lemma-characterize-integral}{{01ON}{source}{}{sourcepointer-4}{}}
\newlabel{properties-lemma-characterize-integral@cref}{{[originallemma][0][]01ON}{[1][0][]source}{}{}{}}
\makeatother
\crefname{originallemma}{original lemma}{original lemmas}
\Crefname{originallemma}{Original lemma}{Original lemmas}
\makeatletter
\newlabel{properties-lemma-characterize-irreducible}{{01OM}{source}{}{sourcepointer-5}{}}
\newlabel{properties-lemma-characterize-irreducible@cref}{{[originallemma][0][]01OM}{[1][0][]source}{}{}{}}
\makeatother
\crefname{originallemma}{original lemma}{original lemmas}
\Crefname{originallemma}{Original lemma}{Original lemmas}
\makeatletter
\newlabel{properties-lemma-characterize-reduced}{{01OL}{source}{}{sourcepointer-6}{}}
\newlabel{properties-lemma-characterize-reduced@cref}{{[originallemma][0][]01OL}{[1][0][]source}{}{}{}}
\makeatother
\crefname{originallemma}{original lemma}{original lemmas}
\Crefname{originallemma}{Original lemma}{Original lemmas}
\makeatletter
\newlabel{schemes-lemma-separated-permanence}{{01KU}{source}{}{sourcepointer-7}{}}
\newlabel{schemes-lemma-separated-permanence@cref}{{[originallemma][0][]01KU}{[1][0][]source}{}{}{}}
\makeatother
\title{Products of varieties over an algebraically closed field are varieties}
\author{The Stacks Project Authors}
\date{stacks-project:varieties.tex:local-snapshot-20261011; GNU FDL 1.2 or later}
\begin{document}
\maketitle
\noindent\textit{Editorial note.} All delivered mathematical text is the author's own, copied from the pinned Stacks Project chapter file varieties.tex; the statement is printed together with the author's own slogan line, and no proof step was rewritten, reordered or removed.\par\medskip
\noindent\textit{Editorial note.} The leading context span is the author's own definition of a variety in the same section; it defines the label definition-variety, which the counted statement uses and which therefore is not listed among the external references.\par\medskip
\noindent\textit{Editorial note.} The chapter's example example-product-not-a-variety (lines 70-77), which shows that the analogous statement over Q is false, is deliberately not printed: the counted lemma is intelligible without it, and the example is not part of the counted proof.\par\medskip
\noindent\textit{Editorial note.} The preamble is the chapter's own retained minimal preamble (sources/stacks-varieties/retained-preamble.tex): a verbatim excerpt of the native Xy-pic setup, of the native theorem declarations and mathematical macros, and of the project's own no-op reference/slogan/history display environments, the last of which the counted statement's slogan needs. The native class line and the xr-hyper externaldocument block are omitted because the assembler supplies a plain amsart wrapper; no macro was invented or redefined.\par\medskip
\noindent\textit{Editorial note.} Locators of the cross-chapter references point at a verbatim excerpt of the project's own tag list, sources/stacks-varieties/referenced-tags.txt, each recorded with its stable Stacks tag.\par\medskip
\noindent\textit{Allowed standard prerequisites (editorial).}\begin{itemize}
\item Hilbert Nullstellensatz and k-rational maximal ideals over an algebraically closed field. Supplies a maximal ideal m inside the nonempty open D(a\_1c\_1) with A/m = k. The reduction to finite affine coverings, the tensor-product set-up and the final contradiction are retained.
\item Characterisations of reduced, irreducible and integral schemes by their open coverings. Used to pass from the integrality of each U\_i x\_k V\_j to integrality of the product; the covering and the product decomposition are part of the counted proof.
\item Permanence of separated and finite type morphisms under base change and composition. Gives separatedness and finite typeness of the structure morphism of the product directly from the corresponding properties of X and Y.
\end{itemize}\medskip
\ifdefined\hypertarget\hypertarget{sourcepointer-1}{}\fi
\noindent\textit{Source reference.} Original theorem 00FV: sources/stacks-varieties/referenced-tags.txt, source lines 4--4. Cross-chapter standard prerequisite cited at proof lines 128-131 to obtain a k-rational maximal ideal in D(a\_1c\_1); the counted proof retains the whole tensor-product argument that consumes it.\par\medskip
\ifdefined\hypertarget\hypertarget{sourcepointer-2}{}\fi
\noindent\textit{Source reference.} Original lemma 01T4: sources/stacks-varieties/referenced-tags.txt, source lines 5--5. Cross-chapter pointer for the finite-typeness of the base change X x\_k Y -> Y (proof lines 96-101); a navigation pointer for a standard permanence statement.\par\medskip
\ifdefined\hypertarget\hypertarget{sourcepointer-3}{}\fi
\noindent\textit{Source reference.} Original lemma 01T3: sources/stacks-varieties/referenced-tags.txt, source lines 6--6. Cross-chapter pointer for finite typeness of the composite X x\_k Y -> Y -> Spec(k) (proof lines 98-101); standard permanence input, not the counted argument.\par\medskip
\ifdefined\hypertarget\hypertarget{sourcepointer-4}{}\fi
\noindent\textit{Source reference.} Original lemma 01ON: sources/stacks-varieties/referenced-tags.txt, source lines 7--7. Cross-chapter characterisation of integrality by open coverings invoked at proof lines 110-112 to lift the affine pieces to the product; the covering argument itself is in the counted proof.\par\medskip
\ifdefined\hypertarget\hypertarget{sourcepointer-5}{}\fi
\noindent\textit{Source reference.} Original lemma 01OM: sources/stacks-varieties/referenced-tags.txt, source lines 8--8. Cross-chapter characterisation of irreducibility by open coverings, cited together with 01OL and 01ON at proof lines 110-112.\par\medskip
\ifdefined\hypertarget\hypertarget{sourcepointer-6}{}\fi
\noindent\textit{Source reference.} Original lemma 01OL: sources/stacks-varieties/referenced-tags.txt, source lines 9--9. Cross-chapter characterisation of reducedness by open coverings, cited together with 01OM and 01ON at proof lines 110-112.\par\medskip
\ifdefined\hypertarget\hypertarget{sourcepointer-7}{}\fi
\noindent\textit{Source reference.} Original lemma 01KU: sources/stacks-varieties/referenced-tags.txt, source lines 10--10. Cross-chapter pointer for separatedness of the composite structure morphism (proof lines 99-103); standard permanence input.\par\medskip
\begin{definition}
\label{definition-variety}
Let $k$ be a field. A {\it variety} is a scheme $X$ over $k$
such that $X$ is integral and the structure morphism
$X \to \Spec(k)$ is separated and of finite type.
\end{definition}

\begin{lemma}
\label{lemma-product-varieties}
\begin{slogan}
Products of varieties are varieties over algebraically closed fields.
\end{slogan}
Let $k$ be an algebraically closed field.
Let $X$, $Y$ be varieties over $k$.
Then $X \times_{\Spec(k)} Y$ is a variety over $k$.
\end{lemma}

\begin{proof}
The morphism $X \times_{\Spec(k)} Y \to \Spec(k)$ is of
finite type and separated because it is the composition of the
morphisms $X \times_{\Spec(k)} Y \to Y \to \Spec(k)$
which are separated and of finite type, see
Morphisms, Lemmas \ref{morphisms-lemma-base-change-finite-type} and
\ref{morphisms-lemma-composition-finite-type}
and
Schemes, Lemma \ref{schemes-lemma-separated-permanence}.
To finish the proof it suffices to show that $X \times_{\Spec(k)} Y$
is integral.
Let $X = \bigcup_{i = 1, \ldots, n} U_i$,
$Y = \bigcup_{j = 1, \ldots, m} V_j$ be finite affine open coverings.
If we can show that each $U_i \times_{\Spec(k)} V_j$ is integral,
then we are done by
Properties, Lemmas \ref{properties-lemma-characterize-reduced},
\ref{properties-lemma-characterize-irreducible}, and
\ref{properties-lemma-characterize-integral}.
This reduces us to the affine case.

\medskip\noindent
The affine case translates into the following algebra statement: Suppose
that $A$, $B$ are integral domains and finitely generated $k$-algebras.
Then $A \otimes_k B$ is an integral domain. To get a contradiction suppose that
$$
(\sum\nolimits_{i = 1, \ldots, n} a_i \otimes b_i)
(\sum\nolimits_{j = 1, \ldots, m} c_j \otimes d_j) = 0
$$
in $A \otimes_k B$ with both factors nonzero in $A \otimes_k B$.
We may assume that $b_1, \ldots, b_n$ are $k$-linearly
independent in $B$, and that $d_1, \ldots, d_m$ are $k$-linearly independent
in $B$. Of course we may also assume that $a_1$ and $c_1$ are nonzero
in $A$. Hence $D(a_1c_1) \subset \Spec(A)$ is nonempty. By the
Hilbert Nullstellensatz
(Algebra, Theorem \ref{algebra-theorem-nullstellensatz})
we can find a maximal ideal $\mathfrak m \subset A$ contained in
$D(a_1c_1)$ and $A/\mathfrak m = k$ as $k$ is algebraically closed.
Denote $\overline{a}_i, \overline{c}_j$ the residue classes of
$a_i, c_j$ in $A/\mathfrak m = k$. The equation above becomes
$$
(\sum\nolimits_{i = 1, \ldots, n} \overline{a}_i b_i)
(\sum\nolimits_{j = 1, \ldots, m} \overline{c}_j d_j) = 0
$$
which is a contradiction with
$\mathfrak m \in D(a_1c_1)$, the linear independence of
$b_1, \ldots, b_n$ and $d_1, \ldots, d_m$, and the fact that $B$ is a domain.
\end{proof}

\end{document}
